\begin{theorem}[Toeplitz 负几何与 strictification counterterm 的正交分裂]\label{thm:conclusion-toeplitz-negative-geometry-strictification-orthogonal-split}
在有限缺陷口径下，设稳定缺陷阶为 \(\kappa\)，并记扫描 Hankel 块
\[
H_{\kappa-1}:=(u_{p+q})_{0\le p,q\le \kappa-1}.
\]
则对充分大 \(N\)，存在可逆矩阵 \(R_N\) 与某个正半定块 \(P_N\succeq 0\)，使
\[
R_N^*T_N(C)R_N=
\begin{pmatrix}
P_N & 0\\
0 & -H_{\kappa-1}^*H_{\kappa-1}
\end{pmatrix}.
\]
从而 Toeplitz 负方向完全由扫描 Hankel Gram
\[
H_{\kappa-1}^*H_{\kappa-1}
\]
决定；更具体地，对任意 \(v\neq 0\)，令
\[
x:=R_N\binom{0}{v},
\]
则有严格负见证
\[
x^*T_N(C)x=-\|H_{\kappa-1}v\|_2^2<0.
\]

另一方面，对任意 \(\eta\in\RR\)，strictification 变换
\[
C^\sharp(w):=C(w)+i\eta
\]
保持核
\[
K(w,w')=\frac{C(w)+C(w')^\ast}{1-ww'^\ast}
\]
以及全部由该核生成的 Toeplitz--PSD 证书、谱读出与负惯性数据不变。故 Toeplitz 机器所能观测到的全部负几何，与参数 \(\eta\) 所携带的 strictification 自由度，分属两个互不相干的层：前者是有限维可层析的签名几何，后者是必须外置的零维账本。
\end{theorem}

\begin{proof}
前半部分正是定理 \ref{thm:conclusion-visible-toeplitz-negative-gram-chart} 的内容：Toeplitz 负方向在充分大截断上由 \(-H_{\kappa-1}^*H_{\kappa-1}\) 精确张成，并给出显式负见证
\[
x^*T_N(C)x=-\|H_{\kappa-1}v\|_2^2.
\]

后半部分则由定理 \ref{thm:conclusion-toeplitz-gauge-blindness-zero-dimensional-ledger-necessity} 给出：变换 \(C\mapsto C+i\eta\) 不改变核 \(K\)，因而不改变任何由 \(K\) 生成的 Toeplitz 截断、PSD 证书与谱读出。故可观测负几何与 strictification 参数 \(\eta\) 必然正交分裂。证毕。
\end{proof}

\begin{corollary}[任何纯 Toeplitz 审计都不可能同时恢复负方向与 strictification 参数]\label{cor:conclusion-no-pure-toeplitz-audit-recovers-negative-directions-and-eta}
不存在任何仅依赖于
\[
\{T_N(C)\}_{N\ge 0}
\]
及由核 \(K\) 生成的全部 Toeplitz--PSD / 谱证书的审计程序，能够一方面完整恢复 Toeplitz 负方向，另一方面又识别 strictification 参数 \(\eta\)。
\end{corollary}

\begin{proof}
由定理 \ref{thm:conclusion-toeplitz-negative-geometry-strictification-orthogonal-split}，Toeplitz 负方向确可由 Hankel Gram 图表精确恢复；但同一定理同时表明 \(C\mapsto C+i\eta\) 在全部此类 Toeplitz 证书上严格不可见。故任何纯 Toeplitz 审计若声称同时恢复二者，便与 \(\eta\) 的不可观测性矛盾。证毕。
\end{proof}

\begin{theorem}[有限缺陷秩与宏观 Toeplitz 审计之间的二尺度分裂]\label{thm:conclusion-finite-defect-rank-vs-macroscopic-toeplitz-audit-scale-splitting}
在有限点状缺陷且无奇异内因子口径下，设稳定缺陷阶为 \(\kappa\)，共动扫描谱样本为 \((u_n)_{n\in\ZZ}\)，Toeplitz 证书矩阵为 \(T_N\)。若存在 \(N_0\) 使得
\[
\operatorname{rank}(H_N)=\kappa
\qquad (N\ge N_0),
\]
则对一切充分大的 \(N\) 成立：
\[
\nu_-(T_N)=\kappa,
\qquad
T_N \sim_{\mathrm{cong}} P_N\oplus\bigl(-H_{\kappa-1}^\ast H_{\kappa-1}\bigr),
\qquad
P_N\succeq 0.
\]
因此，Toeplitz 负方向的全部信息被一个固定的 \(\kappa\times\kappa\) Hankel 块完全承载。并且该负块可由前缀
\[
(u_0,\dots,u_{2\kappa-2})
\]
精确构造，且在一般位置下此长度是最短的。

另一方面，若记固定图表的 Toeplitz--PSD 证书最小资源为 \(N_{\mathrm{Toep}}(T,\delta_0)\)，共动扫描覆盖证书最小资源为 \(N_{\mathrm{scan}}(T,t,\delta_0)\)，则正文给出
\[
N_{\mathrm{Toep}}(T,\delta_0)\gtrsim
\frac{N_{\mathrm{scan}}(T,t,\delta_0)^2}
{\log\!\bigl(1+(t+\delta_0)^2N_{\mathrm{scan}}(T,t,\delta_0)^2\bigr)},
\qquad
N_{\mathrm{scan}}(T,t,\delta_0)\asymp \frac{T}{t+\delta_0}.
\]
故当 \(\kappa\) 有界而 \(T\to\infty\) 时，出现严格的二尺度分裂：内禀缺陷数据是有限维的，而外部 Toeplitz 审计成本却呈近二次增长。
\end{theorem}

\begin{proof}
前半部分正是定理 \ref{thm:conclusion-toeplitz-negative-geometry-strictification-orthogonal-split} 的第一部分：对充分大 \(N\)，存在合同分解
\[
R_N^\ast T_N(C)R_N=
\begin{pmatrix}
P_N&0\\
0&-H_{\kappa-1}^\ast H_{\kappa-1}
\end{pmatrix},
\qquad
P_N\succeq0,
\]
从而
\[
\nu_-(T_N)=\rank(H_{\kappa-1})=\kappa.
\]
又因为
\[
H_{\kappa-1}=(u_{p+q})_{0\le p,q\le \kappa-1},
\]
其全部条目恰由前缀 \((u_0,\dots,u_{2\kappa-2})\) 决定，故负块可由该长度的扫描前缀精确构造。若再删去最后一项，则一般位置下将无法唯一固定 \(H_{\kappa-1}\) 的最外反对角条目；否则便会与最短窗 \(2\kappa\) 的一般位置非唯一性相冲突，故该长度 generically 最短。

后半部分来自正文的 Toeplitz--scan 资源分离律：宏观 Toeplitz 审计若要在高度尺度 \(T\) 上完成 PSD 核验，其最小截断深度至少满足
\[
N_{\mathrm{Toep}}(T,\delta_0)\gtrsim
\frac{N_{\mathrm{scan}}(T,t,\delta_0)^2}
{\log\!\bigl(1+(t+\delta_0)^2N_{\mathrm{scan}}(T,t,\delta_0)^2\bigr)},
\]
而共动扫描覆盖的最小资源满足
\[
N_{\mathrm{scan}}(T,t,\delta_0)\asymp \frac{T}{t+\delta_0}.
\]
将第二式代入第一式便得
\[
N_{\mathrm{Toep}}(T,\delta_0)\gtrsim
\frac{T^2}{(t+\delta_0)^2\log\!\bigl(1+T^2\bigr)},
\]
故当 \(\kappa\) 固定而 \(T\to\infty\) 时，内禀缺陷信息仍由固定大小的 Hankel 块完全承载，而 Toeplitz 审计资源却按近二次尺度增长。证毕。
\end{proof}

\begin{theorem}[Toeplitz 负证书最短编译长度的尾刚性]\label{thm:conclusion-toeplitz-shortest-negative-certificate-tail-rigidity}
在有限缺陷口径下，设 \(C\) 的稳定缺陷阶为
\[
\kappa(C):=\sup_N \nu^{-}\!\bigl(T_N(C)\bigr)\ge 1.
\]
定义其最短负证书编译长度 \(\ell_{\min}(C)\) 为最小整数 \(L\)，满足存在一个统一编译器，仅从扫描指纹喷射
\[
(u_0,\dots,u_L)
\]
出发，便可在充分大 Toeplitz 阶上构造出保证负性的见证向量。则有精确等式
\[
\ell_{\min}(C)=2\kappa(C)-2.
\]
进一步，若两列 Verblunsky 系数 \((\alpha_n)\)、\((\widetilde\alpha_n)\) 自某个 \(N_0\) 起完全相同，则相应对象 \(C,\widetilde C\) 满足
\[
\kappa(C)=\kappa(\widetilde C),
\qquad
\dim \mathcal H_{\mathrm{pp}}(U_\infty)=\dim \mathcal H_{\mathrm{pp}}(\widetilde U_\infty),
\qquad
\ell_{\min}(C)=\ell_{\min}(\widetilde C).
\]
也就是说，最短负证书长度不是前缀现象，而是严格的尾不变量。
\end{theorem}
\begin{proof}
由定理 \ref{thm:conclusion-finite-defect-rank-vs-macroscopic-toeplitz-audit-scale-splitting}，对充分大 \(N\) 有合同分解
\[
T_N(C)\sim_{\mathrm{cong}} P_N\oplus\bigl(-H_{\kappa-1}^\ast H_{\kappa-1}\bigr),
\qquad
P_N\succeq 0,
\]
并且该负块
\[
H_{\kappa-1}=(u_{p+q})_{0\le p,q\le \kappa-1}
\]
完全由前缀
\[
(u_0,\dots,u_{2\kappa-2})
\]
决定。故仅凭长度 \(2\kappa-2\) 的喷射即可统一编译出负见证向量，因而
\[
\ell_{\min}(C)\le 2\kappa(C)-2.
\]
同一定理还明确给出一般位置下的最短性：若删去最后一项 \(u_{2\kappa-2}\)，则无法统一固定最外反对角条目，也就不能同时完成“保证负性”的编译任务。故
\[
\ell_{\min}(C)\ge 2\kappa(C)-2.
\]
两边合并即得
\[
\ell_{\min}(C)=2\kappa(C)-2.
\]

若两列 Verblunsky 系数自某个 \(N_0\) 起完全相同，则由命题 \ref{con:discussion-cmv-tail-invariant}，对应无限维 CMV 对象只差一个有限秩扰动；因此稳定缺陷阶与点谱残差维数都是严格尾不变量：
\[
\kappa(C)=\kappa(\widetilde C),
\qquad
\dim \mathcal H_{\mathrm{pp}}(U_\infty)=\dim \mathcal H_{\mathrm{pp}}(\widetilde U_\infty).
\]
再由已经证明的公式
\[
\ell_{\min}=2\kappa-2,
\]
立得
\[
\ell_{\min}(C)=\ell_{\min}(\widetilde C).
\]
证毕。
\end{proof}

\begin{theorem}[缺陷显影的双参数不可约律]\label{thm:conclusion-toeplitz-defect-revelation-two-parameter-irreducibility}
在结论 \ref{thm:conclusion-toeplitz-shortest-negative-certificate-tail-rigidity} 的有限缺陷设定下，记最浅缺陷深度为
\[
h_{\min}:=\min_{1\le j\le \kappa}(1-|a_j|),
\]
其中 \(\{a_j\}\) 为盘内缺陷节点。定义全显影阶
\[
N_{\mathrm{full}}(C):=\min\{N:\ \nu^{-}(T_M(C))=\kappa(C)\ \forall M\ge N\}.
\]
则存在常数 \(c,C>0\)（只依赖于 \(\kappa\) 与零点分离的粗界）使得
\[
c\,h_{\min}^{-1}\le N_{\mathrm{full}}(C)\le C\,h_{\min}^{-1}.
\]
于是，任何同时完成下列两项任务的统一协议：
\begin{enumerate}
  \item 输出一个保证负性的 Toeplitz 见证；
  \item 证明全部 \(\kappa\) 个负方向已经完全显影；
\end{enumerate}
都必须同时满足两条彼此不可替代的资源下界
\[
L\ge 2\kappa-2,
\qquad
N\ge c\,h_{\min}^{-1},
\]
其中 \(L\) 为扫描喷射长度，\(N\) 为 Toeplitz 截断阶。反之，取
\[
L=2\kappa-2,
\qquad
N\ge C\,h_{\min}^{-1},
\]
即可同时实现这两项目标。

因此，缺陷揭示的最小复杂度并非单参量，而是一个严格分裂为“代数编译长度”与“解析显影深度”的二维资源向量。
\end{theorem}
\begin{proof}
第一条下界
\[
L\ge 2\kappa-2
\]
正是定理 \ref{thm:conclusion-toeplitz-shortest-negative-certificate-tail-rigidity} 的精确结论，且该下界已知可达。

对第二条，定理 \ref{thm:conclusion-visible-complete-defect-audit-dual-scale-threshold} 已给出显影负方向所需的粗尺度阈值：存在常数 \(c,C>0\) 使
\[
N\le c\,h_{\min}^{-1}
\quad\Longrightarrow\quad
\nu^{-}\!\bigl(T_N(C)\bigr)\le \kappa-1,
\]
而当
\[
N\ge C\,h_{\min}^{-1}
\]
时，全部 \(\kappa\) 个负方向都已显影并保持稳定。于是
\[
c\,h_{\min}^{-1}\le N_{\mathrm{full}}(C)\le C\,h_{\min}^{-1}.
\]

现在把两项任务合并考虑。若 \(L<2\kappa-2\)，则统一编译器本身不存在，即使把 \(N\) 取到任意大，也不能完成任务一。若 \(N<c\,h_{\min}^{-1}\)，则 full index 仍不可见，即使 \(L\) 充分长，也不能完成任务二。故这两条资源下界彼此不可替代。

反之，当
\[
L=2\kappa-2,
\qquad
N\ge C\,h_{\min}^{-1},
\]
时，任务一由定理 \ref{thm:conclusion-toeplitz-shortest-negative-certificate-tail-rigidity} 完成，任务二由定理 \ref{thm:conclusion-visible-complete-defect-audit-dual-scale-threshold} 的显影阈值完成，故二者可并行实现。证毕。
\end{proof}
